\documentclass[a4paper]{article}

\usepackage[english]{babel}
\usepackage[T1]{fontenc}
\usepackage[a4paper,top=3cm,bottom=2cm,left=3cm,right=3cm,marginparwidth=1.75cm]{geometry}
\usepackage{amsmath}
\usepackage{amssymb}
\usepackage{graphicx}
\usepackage{booktabs}
\usepackage[colorinlistoftodos]{todonotes}
\usepackage[colorlinks=true, allcolors=blue]{hyperref}
\setlength{\marginparwidth}{2cm}
\setcounter{secnumdepth}{0}
\parindent 0pt
\parskip 1em

\title{Homework 1}
\author{Ng Tze Kean}

\begin{document}

\maketitle

%-----------------------------------------------------------------------
\section{Question 1: A Simple Delegated Monitor Model}
%-----------------------------------------------------------------------

Consider a firm who needs \$100,000 to initiate a project. The payout of
this project is \$200,000 with probability 90\%, \$100,000 with probability
5\%, and \$0 with probability 5\% at the end of two years. There are 1,000
small investors each with \$100 to lend. If the project is initiated, only
the firm manager knows the result of the project, while the lenders cannot
see the results. However, as long as one of the lenders pays a monitor fee
of \$200, the results will be known to everyone.

%-----------------------------------------------------------------------
\subsection{Part (a)}
%-----------------------------------------------------------------------

\textit{Assume each investor knows that there are many other small investors
with \$100 to lend, but each investor lives on an isolated island and cannot
coordinate with other investors; he or she can only communicate with the
firm. Suppose the project is the only investment opportunity for the small
investors. Will the small investors lend directly to the firm? Will the
small investor pay the monitor fee? Why or why not?}

\medskip
\textbf{Step 1: Will investors pay the monitoring fee?}

To determine whether any investor would pay the \$200 monitoring fee, we
analyse the strategic decision using the following Nash Equilibrium payout
matrix:

\begin{center}
\begin{tabular}{lll}
\toprule
\textbf{Your Decision} & \textbf{Others Monitor} & \textbf{Others Don't Monitor} \\
\midrule
\textbf{Monitor}      & Your share $-$ \$200  & Your share $-$ \$200 \\
\textbf{Don't Monitor} & Your share $-$ \$0   & \$0 (firm claims failure) \\
\bottomrule
\end{tabular}
\end{center}

If others are already monitoring, the results are publicly revealed and
paying \$200 is redundant; the dominant strategy is not to monitor. If
others are not monitoring, the firm manager can freely claim the project
failed and pocket the proceeds. An investor's expected share is only
\(0.001 \times (\$200{,}000 \times 0.9 + \$100{,}000 \times 0.05 + \$0 \times
0.05) = \$185\), which is less than the \$200 monitoring cost. Even if one
investor were willing to pay, the free-rider problem ensures no single
investor voluntarily bears the full cost. The unique Nash Equilibrium is
therefore that \textbf{nobody monitors}.

\medskip
\textbf{Step 2: Will investors lend to the firm?}

Given that no monitoring will occur, the firm manager always has the
incentive to claim the project failed and retain all proceeds. A rational
investor anticipates this and expects a return of \$0 on any \$100 lent. Since
keeping the \$100 yields \$100 with certainty while lending yields \$0,
\textbf{no small investor will lend to the firm}. This outcome is a classic
market failure driven by the combination of information asymmetry and
coordination failure: each investor hopes another will pay the monitoring
fee, nobody does, and the project goes unfunded.

%-----------------------------------------------------------------------
\subsection{Part (b)}
%-----------------------------------------------------------------------

\textit{Will the project be funded if the 1,000 small investors can
communicate with each other? If so, explain how they can do it and what is
the expected rate of return on the investment.}

\medskip
Yes, the project will be funded once communication is possible, because
investors can solve both the free-rider problem and the monitoring cost problem
through coordination.

\medskip
\textbf{Option 1 --- Shared Monitoring Cost.}
All 1,000 investors agree to contribute \$0.20 each toward monitoring,
raising the required \$200 in total. Each investor's total outlay becomes
\$100.20. With monitoring guaranteed, the firm manager can no longer
misreport results, and the investment is viable.

\medskip
\textbf{Option 2 --- Delegated Monitor (Financial Intermediary).}
Investors pool their funds and appoint a single intermediary (a bank) that
collects \$100,000 and pays the \$200 monitoring fee on their behalf. The
intermediary acts as a single lender with full monitoring capability,
illustrating the fundamental role of banks as delegated monitors.

\medskip
\textbf{Expected Rate of Return Calculation.}
The gross expected payout of the project is:
\[
\mathbb{E}[\text{Payout}] = \$200{,}000 \times 0.9 + \$100{,}000 \times 0.05
+ \$0 \times 0.05 = \$185{,}000.
\]
The total cost of the investment (principal plus monitoring) is
\(\$100{,}000 + \$200 = \$100{,}200\). The net profit is therefore
\(\$185{,}000 - \$100{,}200 = \$84{,}800\), giving a return over two years of:
\[
\frac{\$84{,}800}{\$100{,}200} \approx 84.63\%.
\]
The annualised return is approximately \((1.8463)^{0.5} - 1 \approx 36.0\%\)
per year. Each investor who contributed \$100.20 can expect to receive
approximately \$185 after two years, a dramatically better outcome than the
\$0 return available without coordination.

\newpage
%-----------------------------------------------------------------------
\section{Question 2: Diversification of Risk through Financial Intermediation}
%-----------------------------------------------------------------------

In an economy there are many identically distributed and independent
projects with a one-year investment horizon. The net rate of return on each
project follows a normal distribution with mean \(\mu = 6\%\) and standard
deviation \(\sigma = 15\%\). Each project requires an initial investment of
\$1,000. Formally, \(r \sim N(\mu, \sigma^{2})\) with \(\mu = 6\%\) and
\(\sigma = 15\%\).

%-----------------------------------------------------------------------
\subsection{Part (a)}
%-----------------------------------------------------------------------

\textit{Suppose a bank invests in only one project. What is the probability
of losing more than half of the initial investment (\$500)? What is the
probability of losing more than 10\% of the initial investment (\$100)?}

\medskip
A loss of more than \$500 corresponds to a return below \(-50\%\). The
standardised $z$-score is:
\[
z = \frac{-0.50 - 0.06}{0.15} = \frac{-0.56}{0.15} \approx -3.733,
\]
so \(P(r < -50\%) = \Phi(-3.733) \approx 0.0094\%\).

A loss of more than \$100 corresponds to a return below \(-10\%\). The
standardised $z$-score is:
\[
z = \frac{-0.10 - 0.06}{0.15} = \frac{-0.16}{0.15} \approx -1.067,
\]
so \(P(r < -10\%) = \Phi(-1.067) \approx 14.31\%\). There is therefore a
non-trivial 14.3\% chance that a single-project bank loses more than 10\%
of its investment.

%-----------------------------------------------------------------------
\subsection{Part (b)}
%-----------------------------------------------------------------------

\textit{Suppose the bank invests in 15 such projects. What is the mean and
standard deviation of the rate of return? What is the probability of losing
more than 10\% of the total initial investment (\$1,500)?}

\medskip
By the Central Limit Theorem, the average return across 15 i.i.d.\ projects
is:
\[
\bar{r}_{15} \sim N\!\left(\mu,\; \frac{\sigma^{2}}{N}\right),
\quad \mu = 6\%,\quad \frac{\sigma}{\sqrt{15}} = \frac{15\%}{\sqrt{15}}
\approx 3.873\%.
\]
A loss of more than \$1,500 on the \$15,000 total investment requires the
average return to fall below \(-10\%\):
\[
z = \frac{-0.10 - 0.06}{0.03873} \approx -4.131,
\]
so \(P(\bar{r}_{15} < -10\%) = \Phi(-4.131) \approx 0.00180\%\). Diversifying
across 15 projects reduces the probability of a 10\% loss from 14.3\% to
just 0.0018\%---a dramatic reduction in risk illustrating the power of
portfolio diversification.

%-----------------------------------------------------------------------
\subsection{Part (c)}
%-----------------------------------------------------------------------

\textit{What is the mean and standard deviation of the rate of return if the
bank invests in $N$ projects? What happens as $N \to \infty$?}

\medskip
For a portfolio of \(N\) independent, identically distributed projects, the
average return satisfies:
\[
\bar{r}_{N} \sim N\!\left(\mu,\; \frac{\sigma^{2}}{N}\right).
\]
The mean remains constant at \(\mu = 6\%\) regardless of \(N\), while the
standard deviation equals \(\sigma/\sqrt{N} = 15\%/\sqrt{N}\), which
decreases monotonically toward zero as \(N\) grows. As \(N \to \infty\),
the Law of Large Numbers guarantees that the portfolio return converges
almost surely to the mean:
\[
\bar{r}_{N} \xrightarrow{\text{a.s.}} \mu = 6\%.
\]
This is the core insight of financial intermediation: a bank holding
infinitely many uncorrelated projects faces zero variance risk, and the
portfolio return becomes a certainty of 6\%.

%-----------------------------------------------------------------------
\subsection{Part (d)}
%-----------------------------------------------------------------------

\textit{Suppose the bank's assets consist of infinitely many such projects
(\(N \to \infty\)), and it offers depositors a one-year certificate of
deposit (CD) with a fixed gross return \(1+r\). What is the maximum
interest rate a fully competitive bank (zero profit) can offer? Ignore
operational costs.}

\medskip
When \(N \to \infty\), the bank earns exactly \(\mu = 6\%\) on its asset
portfolio with no uncertainty. A fully competitive bank earns zero profit by
passing the entire return to depositors:
\[
\text{Maximum CD rate} = \mu = 6\%.
\]
The bank has transformed risky individual loans (each with a 15\% standard
deviation) into a riskless liability---a fixed-return CD---by pooling and
diversifying. This perfectly illustrates how banks create safe secondary
securities (deposits) from risky primary assets (loans), which is a
fundamental function of financial intermediation.

\newpage
%-----------------------------------------------------------------------
\section{Question 3: Bank Run}
%-----------------------------------------------------------------------

Consider a simplified version of the Diamond--Dybvig (1983) model. A bank
collects deposits of 100 from many identical depositors at date~0. At
date~1, a fraction of 20\% of depositors are impatient and must consume
early; the remaining 80\% are patient and can wait until date~2. The bank
invests all funds in a long-term project that yields 1.4 per unit if held to
date~2, or 0.8 per unit if liquidated early at date~1. The demand deposit
contract promises 1.0 per depositor who withdraws at date~1 and 1.2 per
depositor who waits until date~2. The setup is summarised below.

\begin{center}
\begin{tabular}{ll}
\toprule
\textbf{Parameter} & \textbf{Value} \\
\midrule
Total deposits (date 0)              & 100 \\
Impatient depositors (withdraw date 1) & 20\% $\Rightarrow$ 20 depositors \\
Patient depositors (withdraw date 2)   & 80\% $\Rightarrow$ 80 depositors \\
Long-term yield (held to date 2)       & 1.4 per unit \\
Early liquidation yield (date 1)       & 0.8 per unit \\
Deposit promise (date 1)               & 1.0 per depositor \\
Deposit promise (date 2)               & 1.2 per depositor \\
\bottomrule
\end{tabular}
\end{center}

%-----------------------------------------------------------------------
\subsection{Part (a): Solvency under Normal Withdrawal}
%-----------------------------------------------------------------------

\textit{Suppose only the 20 impatient depositors withdraw at date 1. Is the
bank solvent?}

\medskip
To meet the 20 units owed at date~1, the bank must liquidate
\(20 / 0.8 = 25\) units of its investment early, leaving 75 units intact.
At date~2, those 75 units grow to \(75 \times 1.4 = 105\) units, which
comfortably covers the \(80 \times 1.2 = 96\) units owed to patient
depositors:
\[
\text{Date-2 assets} = 105 \;\geq\; 96 = \text{Date-2 obligations.}
\]
The bank is fully solvent and retains a surplus of 9 units. This
demonstrates that a healthy bank faces only a \emph{liquidity} mismatch
between short-term liabilities and long-term assets, not an insolvency
problem.

%-----------------------------------------------------------------------
\subsection{Part (b): Bank Run Equilibrium}
%-----------------------------------------------------------------------

\textit{(i) What happens if patient depositors panic and withdraw at date 1?}

If patient depositors fear that others will run, their dominant strategy
becomes withdrawing early as well. The bank now faces all 100 depositors
demanding 1.0 at date~1, but its assets yield only 0.8 per unit when
liquidated early. Total assets available amount to \(100 \times 0.8 = 80\)
units against obligations of \(100 \times 1.0 = 100\) units, leaving a
shortfall of 20 units. Under a sequential service constraint (first-come,
first-served), the first 80 depositors in line each receive 1.0, while the
last 20 depositors receive nothing and the bank fails:
\[
\text{Assets} = 80 < 100 = \text{Obligations} \;\Rightarrow\; \text{Bank fails.}
\]

\textit{(ii) Why is this a self-fulfilling equilibrium?}

This outcome is a self-fulfilling prophecy because the belief itself causes
the outcome it predicts. If a patient depositor believes others will run, it
is rational to withdraw early and receive 1.0 rather than risk receiving 0
later. Because enough depositors act on this belief, the bank truly fails,
which retroactively justifies the original fear. The model therefore
exhibits two Nash Equilibria: a \emph{good equilibrium} in which patient
depositors wait, the bank remains solvent, and everyone is paid according to
the contract; and a \emph{bad equilibrium} (bank run) in which everyone
withdraws at date~1 and the bank collapses. Neither equilibrium depends on
the bank's fundamental health---the bad equilibrium is driven purely by
coordination failure and panic, which is Diamond and Dybvig's central
insight.

%-----------------------------------------------------------------------
\subsection{Part (c): Policy Interventions}
%-----------------------------------------------------------------------

\textit{(i) How does deposit insurance change the equilibrium?}

Deposit insurance eliminates the bad equilibrium entirely by removing the
incentive to run. With government-guaranteed deposits, a patient depositor
knows they will receive 1.2 at date~2 regardless of what others do.
Withdrawing early yields only 1.0, so waiting is always the dominant
strategy and no rational depositor will run. The key mechanism is that
deposit insurance breaks the strategic complementarity: your optimal action
no longer depends on what you fear others will do. The cost is moral
hazard---banks may take excessive risks knowing deposits are
guaranteed---which is why deposit insurance is typically paired with
prudential regulation and capital requirements.

\medskip
\textit{(ii) Would a lender of last resort (LOLR) eliminate the bad
equilibrium?}

A central bank acting as lender of last resort can also eliminate the bad
equilibrium, but only under specific conditions. It must lend freely at a
penalty rate (to ensure only solvent-but-illiquid banks borrow), against
good collateral (to limit central bank risk), and the bank must be illiquid
rather than insolvent. Crucially, the mere credible announcement of LOLR
support can prevent a run from starting entirely. In this model the bank is
solvent (as shown in Part~a), so a LOLR pledge that the central bank will
provide liquidity if needed makes it irrational for patient depositors to
run---they know the bank can always honour 1.2 at date~2. If the bank were
truly insolvent, however, LOLR support would merely delay failure and create
moral hazard without resolving the underlying problem.

\end{document}
